
Consider a pipe of length \gls{ell}, where \gls{z} $\in [0, \ell]$ denotes position along the pipe. Assume steady-state operation, and let
\[
h : [0,\ell] \to \mathbb{R}_{\geq0}, \quad q : [0,\ell] \to \mathbb{R}_{\geq0}
\]
denote the piezometric head and volumetric discharge (from now regarded as head and flow for brevity). Thus, \gls{hz} and \gls{qz} represent the head and flow at position \gls{z}, which are both assumed strictly positive along the pipe. Suppose \gls{n} leaks occur at locations
\[
0 < z_1 < \dots < z_n < \ell,
\]
with corresponding leak discharges $\tilde q_1,\dots,\tilde q_n$, where \gls{tildeqi} and \gls{zi} correspond to the $i$-th leak. Assuming small leak openings, each leak is modeled as a valve by
\[
\tilde{q}_i = c_i \sqrt{h(z_i)},
\]
where \gls{ci} $>0$ is a constant effective leak coefficient characterizing the $i$-th leak. The leak parameters are then defined by
\begin{equation*}
    \text{\gls{leakparam}} := \begin{bmatrix}
    z_1 & \cdots & z_n & c_1 & \cdots & c_n
\end{bmatrix}^\top \in \mathbb{R}^{2n}.
\end{equation*}
The steady-state head loss along the pipe is governed by
\begin{equation}
\label{eq:task_desc_ufunc}
\frac{dh}{dz} = -U(q(z)),
\end{equation}
where \gls{U} $:[0,\infty)\to[0,\infty)$ is smooth and strictly increasing. For notation purposes, let the head and flow at the inlet and outlet be defined by
\[
\text{\gls{h0}} := h(0), \qquad \text{\gls{q0}} := q(0), \qquad
\text{\gls{hl}} := h(\ell), \qquad \text{\gls{ql}} := q(\ell).
\]
Since flow is conserved between leak locations, $q$ is piecewise constant. Introducing $z_0:=0$ and $z_{n+1}:=\ell$, the flow is given by
\[
q(z)=\text{\gls{qi}}:=
\begin{cases}
q_0, & \text{for } z \in [0, z_1), \quad i = 0, \\[6pt]
q_0 - \sum_{k=1}^{i} \tilde{q}_k, & \text{for } z \in [z_i, z_{i+1}), \quad i = 1, \dots, n,
\end{cases}
\]
where also $q_n=q_{\ell}$ for notation consistency. From \cref{eq:task_desc_ufunc}, it follows that the head is continuous and piecewise linear along the pipe. Let
\[
L_{i+1} := z_{i+1}-z_{i},
\qquad i=0,\dots,n,
\]
denote the lengths of the non-leaking segments, so that
\[
\sum_{i=1}^{n+1} \text{\gls{Li}} = \ell.
\]
The head at the $i$-th leak location is denoted by
\[
\text{\gls{hi}} := h(z_i), \qquad i=1,\dots,n,
\]
and for notation consistency, let $h_{n+1}:=h_{\ell}$. The notation for a pipe is illustrated in \autoref{fig:leak_notation}.\\

\begin{figure}[h!]
    \centering
    \input{figures/tikz_pipe}
    \caption{Notation used for a pipe with $n$ leaks.}
    \label{fig:leak_notation}
\end{figure}

Suppose the head and flow at the pipe's inlet and outlet can be measured at \gls{m} different steady-state operating points. For the $j$-th operating point, let
\begin{equation*}
\text{\gls{y0}} :=\begin{bmatrix}
    h_0^{(j)} \\[0.4em]
    q_0^{(j)}
\end{bmatrix}, \quad
\text{\gls{yl}} := \begin{bmatrix}
    h_{\ell}^{(j)} \\[0.4em]
    q_{\ell}^{(j)}
\end{bmatrix}, \quad
j=1,\dots,m,
\end{equation*}
define the vectors of inlet and outlet measurements. The complete matrices of inlet and outlet measurements are defined by
\[
\text{\gls{M0}}
:=
\begin{bmatrix}
y_0^{(1)} & \cdots & y_0^{(m)}
\end{bmatrix} \in\mathbb{R}^{2\times m},
\qquad
\text{\gls{Ml}}
:=
\begin{bmatrix}
y_{\ell}^{(1)} & \cdots & y_{\ell}^{(m)}
\end{bmatrix} \in\mathbb{R}^{2\times m}.
\]
Assuming \gls{y0} can be controlled, it is later shown that an estimate of the outlet head and flow, \gls{ylhat}, can be uniquely defined by \gls{leakparam} through some function
\begin{equation*}
    \hat{y}_\ell^{(j)} = f(\vartheta,y_0^{(j)}).
\end{equation*}

\begin{definition}{Leak Identification Problem}{leak-identification}
Assume there exists a function $\mathcal{F} : D \to \mathbb{R}^{2\times m}$, where $D\subset \mathbb{R}^{2n}\times\mathbb{R}^{2\times m}$ is the domain containing all allowed pairs (\gls{leakparam}, \gls{M0}). Let \gls{M0} be fixed, and define the function as
\[
\mathcal{F}_{\mathcal{M}_0}(\vartheta) := \begin{bmatrix}
    f(\vartheta,y_0^{(1)}) & \cdots & f(\vartheta,y_0^{(m)})
\end{bmatrix} =
\begin{bmatrix}
\hat{y}_{\ell}^{(1)} & \cdots & \hat{y}_{\ell}^{(m)}
\end{bmatrix}
\]
The leak identification problem is the inverse problem of recovering leak parameters $\vartheta$ from some obtained and consistent measurements \gls{M0} and \gls{Ml}, such that
\[
\mathcal{F}_{\mathcal{M}_0}(\vartheta)=\mathcal{M}_\ell.
\]
\end{definition}
\vspace{3mm}

For notation purposes, let all measurements be defined by
\[
\text{\gls{M}}
:=
\begin{bmatrix}
\mathcal{M}_0 \\[0.4em] 
\mathcal{M}_{\ell}
\end{bmatrix} \in\mathbb{R}^{4\times m} 
\]
and \gls{y} $:=\begin{bmatrix}
    \text{\gls{y0}} \\[0.4em]
    \text{\gls{yl}}
\end{bmatrix} \in \mathbb{R}^4$ is defined as the $j$-th column of \gls{M}.\\

\begin{samepage}
\subsubsection*{The following tasks should be performed:}
\begin{enumerate}
    \item Literature review (see \cite{peralta2023}, \cite{molno2024parallel}, \cite{molno2024conditions} for a starting point). \\

    \item Assuming that the function \gls{U}, the number of leaks \gls{n}, and the pipe length \gls{ell} are known, develop and implement a method for solving the leak identification problem as defined in \autoref{def:leak-identification}.\\

    \item Investigate the maximum number of leaks feasible to identify.\\
    
    \item Perform simulated tests to investigate the performance of the method in the idealized setup (\gls{U}, \gls{n}, and \gls{ell} are known).\\

    \item Investigate how the method performs when parameters are uncertain, or noise is introduced to measurements.\\
\end{enumerate}
\end{samepage}